\chapter{Notazione}
\label{ch:notazione}

\section{Indici e insiemi}

\begin{center}
\begin{tabularx}{\textwidth}{@{}l l X@{}}
\toprule
\textbf{Simbolo} & \textbf{Dominio} & \textbf{Significato} \\
\midrule
$t$        & $0,\dots,T-1$ & indice dello slot temporale \\
$T$        & $\mathbb{N}$  & numero di slot in un giorno (96 con slot da 15 minuti) \\
$i$        & $1,\dots,n$   & indice degli \emph{event device} (elettrodomestici a ciclo) \\
$j$        & $1,\dots,m$   & indice dei device \emph{always-on} (frigoriferi, congelatore) \\
$k$        & $\{0,1,2\}$   & indice del livello di potenza (basso, medio, alto) \\
$K$        & $=3$          & numero di livelli di potenza \\
$\Vset$    & $\subseteq \{0,\dots,T-1\}$ & slot con segnale valido (non mancante) \\
\bottomrule
\end{tabularx}
\end{center}

Tutte le variabili associate a uno slot non valido ($t \notin \Vset$) sono fissate a zero
tramite i limiti superiori delle colonne. In questo modo i buchi di misura non generano
attivazioni fittizie né contribuiscono all'errore di ricostruzione.

\section{Dati del problema}

\begin{center}
\begin{tabularx}{\textwidth}{@{}l l X@{}}
\toprule
\textbf{Simbolo} & \textbf{Unità} & \textbf{Significato} \\
\midrule
$y_t$        & W      & segnale aggregato misurato allo slot $t$ \\
$p_i$        & W      & potenza tipica del device $i$ (\code{p\_typical\_w}) \\
$p_{i,k}$    & W      & potenza del device $i$ al livello $k$ \\
$\delta$     & ---    & scarto frazionario dei livelli (\code{power\_level\_variation}, 0.15) \\
$a_i$        & slot   & durata minima di accensione, $\lceil \code{dur\_min\_min} / g \rceil$ \\
$b_i$        & slot   & durata massima di accensione consecutiva \\
$g$          & minuti & ampiezza dello slot (\code{granularity\_min}, default 15) \\
$W_i$        & ore    & ampiezza della fascia oraria dichiarata (24 se non dichiarata) \\
$E_i$        & slot   & monte ore giornaliero atteso dal questionario \\
$r_i$        & ---    & accensioni ancora disponibili nella settimana \\
$s_i$        & slot   & durata tipica di un ciclo, $\code{dur\_typical\_min} / g$ \\
$f_i(t)$     & ---    & fattore temporale del device $i$ allo slot $t$ (cap.~\ref{ch:fattore}) \\
\bottomrule
\end{tabularx}
\end{center}

I tre livelli di potenza sono costruiti simmetricamente attorno al valore tipico:
\begin{equation}
  p_{i,0} = p_i (1-\delta), \qquad
  p_{i,1} = p_i, \qquad
  p_{i,2} = p_i (1+\delta).
  \label{eq:livelli}
\end{equation}
Con $\delta = 0.15$ ogni apparecchio può assorbire il $15\%$ in più o in meno rispetto al valore
nominale. Questo assorbe la variabilità reale dei consumi (tensione di rete, carico, usura) senza
introdurre variabili continue di potenza, che renderebbero il problema molto più difficile.

\section{Variabili di decisione}

\begin{center}
\begin{tabularx}{\textwidth}{@{}l l X@{}}
\toprule
\textbf{Simbolo} & \textbf{Tipo} & \textbf{Significato} \\
\midrule
$x_{i,t}$        & binaria    & il device $i$ è acceso allo slot $t$ \\
$z_{i,t,k}$      & binaria    & il device $i$ opera al livello $k$ allo slot $t$ \\
$z^{\text{ao}}_{j,t,k}$ & binaria & il device always-on $j$ opera al livello $k$ allo slot $t$ \\
$u_{i,t}$        & binaria    & fronte di salita: il device $i$ si accende in $t$ \\
$\dw_{i,t}$      & binaria    & fronte di discesa: il device $i$ si spegne in $t$ \\
$\ep_t$          & continua $\ge 0$ & sovrastima della ricostruzione allo slot $t$ \\
$\en_t$          & continua $\ge 0$ & sottostima della ricostruzione allo slot $t$ \\
$s^{-}_{i,t}$    & continua $\ge 0$ & slack di durata minima (ciclo troppo corto) \\
$s^{+}_{i,t}$    & continua $\ge 0$ & slack di durata massima (ciclo troppo lungo) \\
$d^{+}_i, d^{-}_i$ & continua $\ge 0$ & scostamento dal monte ore giornaliero atteso \\
$e_i$            & continua $\ge 0$ & accensioni eccedenti la quota settimanale \\
\bottomrule
\end{tabularx}
\end{center}

\begin{notebox}{Corrispondenza con il codice}
Le variabili $\ep_t$ ed $\en_t$ compaiono nel codice come \code{ep} ed \code{en}; $\dw_{i,t}$ come
\code{dw}. Nei modelli che distinguono i due gruppi di apparecchi le variabili degli event device
portano il suffisso \code{\_ev} e quelle degli always-on il suffisso \code{\_ao}.
\end{notebox}

\section{Coefficienti di penalità}

\begin{center}
\begin{tabularx}{\textwidth}{@{}l l X@{}}
\toprule
\textbf{Simbolo} & \textbf{Nome nel codice} & \textbf{Penalizza} \\
\midrule
$\lambda_w$     & \code{time\_window\_penalty}     & accensione fuori dalla fascia oraria dichiarata \\
$\lambda_a$     & \code{activation\_penalty}       & ogni singola accensione \\
$\lambda_d$     & \code{duration\_penalty\_block}  & durata del ciclo fuori dai limiti dichiarati \\
$\lambda_D$     & \code{duration\_penalty\_daily}  & monte ore giornaliero diverso dall'atteso \\
$\lambda_{q,i}$ & derivato                          & accensioni oltre la quota settimanale \\
$\lambda_s$     & \code{season\_penalty}           & uso fuori dai mesi dichiarati \\
$w$             & \code{window\_vagueness\_weight} & genericità della fascia oraria dichiarata \\
\bottomrule
\end{tabularx}
\end{center}

Tutti i coefficienti sono \textbf{adimensionali} e vengono moltiplicati per $p_i$ prima di entrare
nell'obiettivo, come discusso nel capitolo~\ref{ch:perche-l1}.
